[Geht man von Punkt 1 oder 2 aus (Band IV, S.~458), kann
man nun mithilfe der Formeln
\[
x_\nu-x_\mu=\mathrm{s}_{\mu.\nu}\cos\tau_{\mu.\nu}
\qquad
y_\nu-y_\mu=\mathrm{s}_{\mu.\nu}\sin\tau_{\mu.\nu},
\]
in denen \(\tau_{\mu.\nu}\) den Azimut der Seite \(\mathrm{s}_{\mu.\nu}\) bezeichnet, die ebenen Koordinaten der
Punkte des Kranzsystems berechnen. Man erhält:]
\begin{center}
\begin{tabular}{|c|r|r|}
\hline
\shortstack{\(\Bigl[\)Dreiecks-\\punkt}&\multicolumn{1}{c|}{\shortstack{\(x\)\\(Meter)}}&\multicolumn{1}{c|}{\shortstack{\(y\quad\Bigr]\)\\(Meter)}}\\
\hline
21&\(-173074{,}708\)&\(+76350{,}926\)\\
1&\(-196973{,}309\)&\(+44130{,}578\)\\
2&\(-163594{,}034\)&\(+44884{,}912\)\\
3&\(-138675{,}054\)&\(+63178{,}019\)\\
\end{tabular}
\end{center}

\clearpage
\noindent336\hfill\textbf{NACHLASS.}
\begin{center}
\begin{tabular}{|c|r|r|}
\hline
\shortstack{\(\Bigl[\)Dreiecks-\\punkt}&\multicolumn{1}{c|}{\shortstack{\(x\)\\(Meter)}}&\multicolumn{1}{c|}{\shortstack{\(y\quad\Bigr]\)\\(Meter)}}\\
\hline
4&\(-142251{,}677\)&\(+87900{,}386\)\\
5&\(-117302{,}389\)&\(+73520{,}724\)\\
6&\(-115364{,}146\)&\(+117156{,}400\)\\
7&\(-82615{,}970\)&\(+99921{,}647\)\\
8&\(-73684{,}811\)&\(+129246{,}257\)\\
9&\(-114711{,}937\)&\(+148300{,}127\)\\
10&\(-147940{,}705\)&\(+128490{,}294\)\\
11&\(-153046{,}113\)&\(+162443{,}404\)\\
12&\(-194283{,}391\)&\(+134463{,}972\)\\
13&\(-192087{,}067\)&\(+166477{,}506\)\\
14&\(-218810{,}103\)&\(+163540{,}108\)\\
15&\(-229342{,}324\)&\(+135140{,}342\)\\
16&\(-209472{,}265\)&\(+120150{,}094\)\\
17&\(-232173{,}634\)&\(+108214{,}579\)\\
18&\(-227661{,}782\)&\(+89453{,}677\)\\
19&\(-193865{,}475\)&\(+81844{,}596\)\\
20&\(-209919{,}630\)&\(+63160{,}451\)\\
21&\(-173075{,}289\)&\(+76350{,}558\)\\
1&\(-196973{,}026\)&\(+44131{,}372\)\\[6pt]
&\multicolumn{2}{c|}{Unterschiede}\\[3pt]
21&\(-0{,}581\)&\(-0{,}368\)\\
1&\(+0{,}283\)&\(+0{,}794\)\\
\end{tabular}
\end{center}
